\section{Hàm mất mát và hậu xử lý}
Trong YOLO11, huấn luyện tối ưu đồng thời định vị hộp, phân loại và phân phối khoảng cách biên. Cấu hình E3/E4 đặt hệ số box 7,5; class 0,5 và DFL 1,5. Các hệ số xác định tỷ lệ đóng góp vào gradient, không phải mức quan trọng của lớp trong ứng dụng. Mosaic được bật, biến đổi HSV và lật ngang tạo thêm đa dạng; mixup, cutmix và copy-paste tắt. Hai run dùng cùng thiết lập nên khác biệt n/s được khảo sát trên nền tương đối đồng nhất.

Faster R-CNN trả các loss phân loại, hồi quy hộp, objectness RPN và hồi quy proposal. E6 chuyển nhãn YOLO chuẩn hóa thành hộp pixel \texttt{xyxy}, đặt \texttt{iscrowd=0}, và cộng một vào class ID vì Torchvision dành 0 cho background. Ảnh qua transform với cạnh ngắn 640, cạnh dài tối đa 960. Đây không phải resize vuông như YOLO, nên so sánh phản ánh pipeline hoàn chỉnh chứ không cô lập tuyệt đối kiến trúc.

Sau forward, hộp score thấp bị lọc và NMS loại dự đoán chồng lặp. NMS quá thấp có thể xóa hai vật thể gần nhau, quá cao có thể giữ nhiều hộp trùng. Evaluator thu prediction YOLO với NMS 0,7; Torchvision có hậu xử lý tích hợp. Đây là một nguyên nhân khiến metric log framework và evaluator chung không trùng hoàn toàn.

\section{Lựa chọn ngưỡng bằng F2}
Prediction được thu một lần tại confidence 0,01, tối đa 300 detection/ảnh, rồi cache được quét qua 95 ngưỡng từ 0,01 đến 0,95. Một prediction là TP nếu cùng lớp, IoU ít nhất 0,5 và ghép với một GT chưa dùng. Khi hai ngưỡng cùng F2, quy tắc ưu tiên Precision rồi ngưỡng cao hơn. Quy tắc tie ghi trước giúp lựa chọn tái lập.

F2 đặt Recall nặng gấp bốn lần Precision trong mẫu số, phù hợp mục tiêu giảm bỏ sót. Nó không làm mọi FN có cùng mức nguy hiểm: bỏ sót bậc thang gần người dùng khác bỏ sót tòa nhà xa, nhưng dữ liệu không có khoảng cách hoặc mức rủi ro. F2 vì vậy là tiêu chí kỹ thuật cho detector, không phải thước đo an toàn di chuyển.

Ngưỡng 0,22/0,19/0,59 khác nhau vì phân phối score của kiến trúc khác nhau. Không thể nói E6 “tự tin hơn” chỉ vì threshold cao. Score chưa được hiệu chuẩn xác suất; 0,8 không có nghĩa đúng 80 phần trăm. Muốn so độ tin cậy cần reliability diagram, expected calibration error hoặc phương pháp calibration phù hợp cho detection.

\section{Học chuyển giao và checkpoint}
Ba mô hình dùng COCO pretrained. E3/E4 nạp \texttt{yolo11n.pt}/\texttt{yolo11s.pt}; E6 nạp COCO V2 mặc định rồi thay predictor bằng 26 đầu ra. Head mới học 25 lớp ROD và background, còn backbone/FPN bắt đầu từ đặc trưng hình ảnh tổng quát.

Ultralytics lưu best theo fitness ưu tiên mAP50--95. E6 lưu best khi mAP50--95 VALID từ TorchMetrics tăng. Các metric này phục vụ huấn luyện nhưng không đặt chung làm bảng xếp hạng vì code chấm khác nhau. Sau huấn luyện, ba best checkpoint được chấm lại bởi cùng evaluator. Recovery checkpoint giữ optimizer, scheduler, scaler và RNG để resume; best weights đã strip optimizer chỉ dùng suy luận.

\section{Nguồn sai số và độ bất định}
Sai số quan sát gồm sai số mô hình, nhãn, threshold, hậu xử lý và biến thiên theo seed. Một con số mAP bốn chữ số thập phân không có bốn chữ số chắc chắn. Báo cáo chưa bootstrap khoảng tin cậy và chưa chạy nhiều seed, nên chênh lệch rất nhỏ theo lớp chỉ dùng gợi ý kiểm tra.

Trường hợp khó dự kiến gồm vật nhỏ, che khuất, ánh sáng yếu, nền giống vật thể và lớp gần nghĩa như Bike/Motorcycle. Đây là giả thuyết từ đặc điểm detection. Muốn xác nhận phải ghép prediction với ảnh, phân nhóm lỗi theo kích thước/điều kiện và duyệt có quy tắc; bảng FN tự nó không chứng minh nguyên nhân.
